\respcap{cap:bfs}
\resp{bl-arv} Fixe uma raiz candidata $r$ (tente cada vértice). \emph{Condição 1}: $\operatorname{dist}_T(r,v)=\operatorname{dist}_G(r,v)$ para todo $v$ (calcule por busca em largura em $G$); ela é necessária, porque o nível na árvore de largura é a distância. \emph{Condição 2 (ordem)}: numa busca em largura, os vértices saem da fila nível a nível; dentro de um nível, irmãos ficam juntos, e os blocos de irmãos seguem a ordem dos pais; e cada $w$ é filho do \emph{primeiro} vizinho de $w$ do nível anterior a sair da fila. Assim, a única liberdade da busca é a ordem dos filhos de cada vértice, e a condição 2 se traduz em restrições ``$x$ antes de $y$'': para cada $w$ e cada outro vizinho $v$ de $w$ no nível anterior, $pai(w)$ deve vir antes de $v$. Se $x$ e $y$ são irmãos, é uma restrição sobre a ordem dos filhos do pai deles; se não são, ela equivale a ``$pai(x)$ antes de $pai(y)$'', e sobe-se até chegar a dois irmãos. $T$ é árvore de largura com raiz $r$ se e somente se, para cada vértice, as restrições sobre a ordem dos seus filhos não formam ciclo, o que se verifica com uma ordenação topológica (Capítulo~\ref{cap:dfsdig}).

\resp{bl-pilha} Com pilha, o vértice consultado é sempre o último inserido que ainda não foi esgotado: a busca avança a partir do vértice mais recente e só volta quando ele não tem vizinhos desmarcados, que é a busca em profundidade; os pais formam a árvore de profundidade. Com fila, o vértice consultado é o mais antigo: todos os seus vizinhos desmarcados são inseridos antes que qualquer outro vértice seja examinado, ou seja, ``esgota-se o pai antes dos filhos'', que é a busca em largura.

\resp{bl-apl} (a) Busca em largura com raiz $x$: $\operatorname{dist}(x,v)=nivel(v)$, e o caminho na árvore é mínimo. (b) Modele cada casa livre como um vértice, ligando casas vizinhas livres; faça a busca a partir da entrada; se a saída for alcançada, o caminho na árvore (subindo pelos pais) é o menor, com $nivel(\text{saída})$ passos. (c) O grafo é bipartido se e somente se a floresta de largura não tem arestas-irmão nem primo, que fecham ciclos ímpares; a 2-coloração é dada pela paridade do nível.

\resp{bl-sim} Ordem de visita: $s,r,w,v,t,x,u,y$. Níveis: $s$ no nível 0; $r,w$ no 1; $v,t,x$ no 2; $u,y$ no 3. Pais: $r,w\leftarrow s$; $v\leftarrow r$; $t,x\leftarrow w$; $u\leftarrow t$; $y\leftarrow x$. Arestas que não são de árvore: $tx$ é irmão (mesmo nível e mesmo pai $w$); $xu$ é tio ($u$ já tinha sido alcançado por $t$, um nível abaixo de $x$); $uy$ é primo (mesmo nível, pais $t$ e $x$). O grafo não é bipartido: o triângulo $w,t,x$ é um ciclo ímpar.

\resp{bl-diam} O vértice $a$ mais distante de qualquer vértice é extremo de algum caminho de comprimento máximo; se não fosse, trocando trechos de caminhos obter-se-ia um caminho mais longo que o diâmetro. Então o vértice mais distante de $a$ dá o diâmetro. No exemplo, a partir do vértice 1, o mais distante é o 7 (distância 3); a partir de 7, o mais distante é o 6, à distância \textbf{5} (caminho $7,5,2,1,3,6$).

\resp{bl-num} Faça $cnt(raiz)=1$. Ao processar uma aresta $vw$ com $nivel(w)=nivel(v)+1$ (aresta-pai ou aresta-tio), faça $cnt(w):=cnt(w)+cnt(v)$. Como todos os vértices do nível $\ell$ são processados antes dos do nível $\ell+1$, $cnt(v)$ já está completo quando $v$ é processado. Custo $O(n+m)$.

\resp{bl-prova} Os vértices na fila pertencem sempre a no máximo dois níveis consecutivos, em ordem não decrescente de nível; por indução, retira-se um vértice de nível $\ell$ e entram no fim vértices de nível $\ell+1$. Seja $vw$ uma aresta, com $v$ processado antes de $w$. Se $w$ ainda não havia sido alcançado, vira filho de $v$ e $nivel(w)=nivel(v)+1$. Se já estava na fila, seu nível é $nivel(v)$ ou $nivel(v)+1$. Ele não pode já ter saído da fila, pois então $w$ teria sido processado antes de $v$. Logo $|nivel(v)-nivel(w)|\le1$. Quanto à distância: o caminho na árvore tem $nivel(v)$ arestas, então $\operatorname{dist}\le nivel(v)$; e qualquer caminho da raiz a $v$ com $k$ arestas muda o nível em no máximo 1 por aresta, então $nivel(v)\le k$.

\resp{bl-grade} Numerando linhas e colunas a partir de 0, o caminho mínimo é $(0,0)\to(0,1)\to(0,2)\to(1,2)\to(2,2)\to(2,3)\to(3,3)$, com \textbf{6} passos. Distâncias: $(0,1)$: 1; $(0,2)$: 2; $(1,2)$: 3; $(2,2)$: 4; $(2,1)$ e $(2,3)$: 5; $(2,0)$, $(3,1)$ e $(3,3)$: 6.

\resp{bl-dig} O algoritmo é o mesmo, percorrendo $N_{out}(v)$; ele calcula a distância (em arestas) da raiz a cada vértice alcançável. Uma aresta $v\to w$ sempre satisfaz $nivel(w)\le nivel(v)+1$, mas pode ``voltar'' vários níveis ($nivel(w)$ muito menor que $nivel(v)$), pois a desigualdade só vale num sentido.

\resp{bl-cam} \textbf{Se} $L(y)=0$, não há caminho. \textbf{Senão}, suba de $y$ por $pai(\cdot)$ até $x$, empilhando os vértices, e depois desempilhe-os, imprimindo de $x$ para $y$ (ou insira cada vértice à esquerda numa lista). Custo $O(nivel(y))$.

\resp{bl-custo} Cada vértice entra na fila e é processado uma vez. Com listas, processar $v$ custa $O(1+\text{grau}(v))$, e a soma dos graus é $2m$: $O(n+m)$. Com matriz, processar $v$ exige percorrer a linha inteira, $O(n)$: $O(n^2)$. As duas coincidem assintoticamente em grafos densos, com $m=\Theta(n^2)$.

\resp{bl-01} Mantenha uma fila de duas pontas com os vértices e rótulos de distância $d$. Ao retirar $v$ do início, para cada aresta $vw$ de peso $p$: se $d(v)+p<d(w)$, atualize $d(w)$ e insira $w$ no \emph{início} se $p=0$, ou no \emph{fim} se $p=1$. A fila mantém a propriedade de conter no máximo dois valores de distância consecutivos, com os menores no início, como a fila de prioridades do algoritmo de Dijkstra; cada vértice é processado um número constante de vezes, e o custo é $O(n+m)$.

\resp{bl-multi} Coloque todas as saídas na fila no início, com nível 0, como se fossem filhas de uma raiz virtual, e execute a busca em largura normalmente. O nível final de cada local é a distância até a saída mais próxima. Custo $O(n+m)$.

\resp{bl-bip} No ciclo $1,2,3,4,5$, a partir de 1: os vértices 2 e 5 ficam no nível 1, e 3 e 4 no nível 2; a aresta $34$ liga dois vértices do nível 2 com pais diferentes (2 e 5): é uma aresta-primo, e o ciclo não é bipartido. No ciclo de 6 vértices, a partir de 1: 2 e 6 no nível 1, 3 e 5 no nível 2 e 4 no nível 3, alcançado por 3; a aresta $54$ é uma aresta-tio, e não há arestas-irmão nem primo: o ciclo é bipartido.

\resp{bl-comp} Use o laço externo da busca em largura com um contador $c$: cada vez que um vértice ainda não alcançado inicia uma nova busca, incremente $c$ e rotule com $c$ todos os vértices alcançados nessa busca. Ao final, $c$ é o número de componentes. Custo $O(n+m)$.

\resp{bl-cavalo} Cada casa é um vértice, ligado às até 8 casas alcançáveis por um movimento de cavalo. Uma busca em largura a partir de a1 dá distância \textbf{6} até h8.

\resp{bl-jarros} Os estados são pares $(a,b)$, com $0\le a\le3$ e $0\le b\le5$, ligados pelas operações possíveis. Uma busca em largura a partir de $(0,0)$ encontra 4 litros com \textbf{6} operações: encher o de 5, despejar no de 3 (sobram 2 no de 5), esvaziar o de 3, despejar os 2 litros no de 3, encher o de 5 e despejar no de 3 até enchê-lo, o que deixa 4 litros no de 5.

\resp{bl-impar} Numa aresta-irmão $vw$, $v$ e $w$ têm o mesmo pai $u$: o ciclo $u,v,w$ tem 3 arestas. Numa aresta-primo $vw$, $v$ e $w$ estão no mesmo nível $\ell$; os caminhos na árvore de $v$ e de $w$ até o seu ancestral comum mais próximo $u$, no nível $j$, têm $\ell-j$ arestas cada; com a aresta $vw$, o ciclo tem $2(\ell-j)+1$ arestas, um número ímpar. Numa aresta-tio, os níveis dos extremos diferem de 1, e o ciclo formado tem comprimento par.

\resp{bl-girth} Faça uma busca em largura a partir de cada vértice $s$. Durante a busca, toda aresta $vw$ que não é de árvore fecha, com os caminhos da árvore até $s$, um passeio fechado de comprimento $nivel(v)+nivel(w)+1$, que contém um ciclo de comprimento no máximo esse valor. O menor valor dessa expressão, sobre todas as raízes e arestas, é exatamente a cintura: quando $s$ pertence a um ciclo mínimo, a aresta desse ciclo ``oposta'' a $s$ dá o seu comprimento exato. Cada busca custa $O(n+m)$, e o total é $O(n(n+m))$.

\resp{bl-k33} A partir de $a$: $x,y,z$ ficam no nível 1 (filhos de $a$) e $b,c$ no nível 2 (filhos de $x$, o primeiro a sair da fila). As arestas $yb$, $yc$, $zb$ e $zc$ são arestas-tio. Não há arestas-irmão nem primo, coerente com o fato de $K_{3,3}$ ser bipartido.
